Residual-quantization (RQ) generative retrievers assign every candidate a short discrete semantic ID and decode it directly, replacing nearest-neighbor search; codebook collapse -- a few codes absorbing most assignments -- is a known failure mode, and balance regularization (an auxiliary loss pushing code usage toward uniform) is the standard remedy. But does more balanced code usage actually improve retrieval? We test this with a controlled ablation of regularization strength on GENIUS, a state-of-the-art RQ-based generative multimodal retriever, across MSCOCO and VisualNews (broad-domain, bidirectional text$\leftrightarrow$image retrieval) and FashionIQ (fine-grained composed retrieval). We find (1) balance regularization does not consistently improve retrieval: it raises MSCOCO text-to-image Recall@1 significantly ($+12.3\%$ relative, $p<10^{-6}$, verified across 5 seeds and 3 independently retrained tokenizers) but hurts FashionIQ and reverses sign on VisualNews ($-53\%$); (2) its effect differs sharply by direction on identical checkpoints -- the same regularization collapses image-to-text Recall from $10.3\%$ to near zero, and a natural per-modality fix rescues neither direction; (3) both effects are already present at the tokenizer level, before decoder training, via a decoder-free probe, and do not transfer consistently across datasets. Absolute Recall in our setup trails published GENIUS numbers, reflecting a matched from-scratch pipeline rather than a state-of-the-art reproduction; our claims concern regularization's relative effect under this fixed pipeline, not absolute performance.
